% ============================================================
% 03_methodology.tex — Methodology (Macedonian)
% ============================================================

\section{Преглед на системската архитектура}

Системот се состои од три примарни слоја кои работат сукцесивно пред секоја натпреварувачка недела, поддржани од intel pipeline кој собира информации во реално време.

\textbf{Слој 1 — Предикција}: Четири LightGBM модели (еден по позиција) предвидуваат очекувани поени за секој играч за следната натпреварувачка недела, користејќи историски и тековни сезонски карактеристики.

\textbf{Слој 2 — Оптимизација}: ILP решавач (PuLP / CBC) ги зема предвидените поени и решава проблем за избор на 15 играчи кои го максимизираат вкупниот предвиден збир на поени, почитувајќи ги сите FPL ограничувања.

\textbf{Слој 3 — Интелигенција}: Комбинација на real-time intel сигнали (достапност, ротација) и LLM препораки (Gemini) ги коригираат предикциите пред оптимизацијата, а Claude API генерира нарациски објаснувања по симулацијата.

Текот на пипелајнот за секоја натпреварувачка недела е следниот:
\begin{enumerate}
  \item Intel 01: Снимање на FPL API (повреди, цени, сопственост).
  \item Intel 02: Скрапирање на прес-конференции (Fantasy Football Scout).
  \item Intel 03: Пресметка на достапност (0--100) по играч.
  \item Intel 04: Пресметка на ризик за ротација (0--100) по играч.
  \item Intel 05: LLM препораки (Gemini 2.5 Flash).
  \item Intel 06: Примена на казни на предикциите $\rightarrow$ скорегирани предикции.
  \item LightGBM: Предикција на поени по играч.
  \item ILP: Избор на оптимален тим од 15 играчи.
  \item Intel 07: Анализа на бенч зголемување (lookahead).
  \item Stage 9: Claude API генерира нарации за одлуките.
\end{enumerate}

\section{Датасет и инженеринг на карактеристики}

\subsection{Историски податоци}

Основата на тренирачкиот датасет ги сочинуваат историски FPL податоци за играчи-по-натпреварувачка недела достапни преку репозиторот на Vaastav \cite{vaastav2024fpl}, кој ги агрегира официјалните FPL API податоци за секоја завршена сезона. Датасетот опфаќа шест сезони: 2019--20, 2020--21, 2021--22, 2022--23, 2023--24 и 2024--25, со вкупно \textbf{51\,044 редови} по примена на филтерот за минимум 30 минути на игра.

Податоците се поделени на четири позициски датотеки:

\begin{table}[h]
\centering
\caption{Распределба на тренирачкиот датасет по позиција}
\label{tab:dataset_split}
\begin{tabular}{lrr}
\toprule
\textbf{Позиција} & \textbf{Редови} & \textbf{Колони} \\
\midrule
Голман (GK)       & 4\,421  & 69 \\
Одбранбен (DEF)   & 18\,828 & 67 \\
Среден (MID)      & 22\,132 & 65 \\
Напаѓач (FWD)     & 5\,663  & 65 \\
\midrule
\textbf{Вкупно}   & \textbf{51\,044} & — \\
\bottomrule
\end{tabular}
\end{table}

Секој ред претставува перформансата на еден играч во една натпреварувачка недела. Целната колона е \texttt{total\_points} — вистинскиот FPL поен на играчот за тоа коло. Важно е да се напомене дека 2025--26 е \textbf{слепа тест} сезона: ниеден податок од тековната сезона не бил инкорпориран во тренирачкиот сет, со цел обезбедување на строга out-of-sample евалуација.

\subsection{Карактеристики за тимски форм}

Тимскиот форм е изведен од Understat, веб-страница која ги снима очекуваните голови (xG) и пропуштени очекувани голови (xGA) за секој натпревар во Премиер Лигата. Бидејќи Understat повеќе не ги вградува податоците директно во HTML-от на страницата, нашиот систем ги повикува JSON API крајните точки директно.

За секој тим и за секој натпревар, се пресметуваат тркалачки прозорци со должина 5 натпревари (последни 5 натпревари):

\begin{equation}
\text{team\_xG\_last5}_t = \frac{1}{5} \sum_{i=t-5}^{t-1} \text{xG}_i
\end{equation}

каде $\text{xG}_i$ е очекуваниот број голови на тимот во натпреварот $i$. Слично се пресметуваат \texttt{team\_xGA\_last5}, \texttt{team\_attacking\_strength}, \texttt{team\_defensive\_strength}, и \texttt{team\_cs\_probability} (веројатност за чист лист изведена од xGA).

\subsection{Карактеристики на нови трансфери}

За нови играчи кои нема историски FPL podatoci (дебитанти во Премиер Лига), системот собира податоци од нивните претходни лиги преку FBref. Покриени се девет европски лиги со соодветни множители кои ја одразуваат конкурентноста на секоја лига во однос на Премиер Лигата:

\begin{table}[h]
\centering
\caption{Множители на лигата за нормализација на статистиките на нови трансфери}
\label{tab:league_multipliers}
\begin{tabular}{lc}
\toprule
\textbf{Лига} & \textbf{Множител} \\
\midrule
La Liga (Шпанија)          & 0.92 \\
Serie A (Италија)           & 0.88 \\
Bundesliga (Германија)      & 0.89 \\
Ligue 1 (Франција)          & 0.82 \\
Championship (Англија)      & 0.72 \\
Primeira Liga (Португалија) & 0.78 \\
Eredivisie (Холандија)      & 0.75 \\
Belgian Pro League          & 0.74 \\
Scottish Premiership        & 0.65 \\
\bottomrule
\end{tabular}
\end{table}

Статистиките (голови, асистенции, минути, интерцепции, победени контакти) се нормализираат на 90 минути и се множат со соодветниот множител пред интегрирање во тренирачкиот сет.

\subsection{Live FPL API податоци}

Пред секоја натпреварувачка недела, системот ги превзема следните live податоци директно од официјалниот FPL API:

\begin{itemize}
  \item \texttt{value} — тековната цена на играчот во милиони фунти.
  \item \texttt{selected\_by\_percent} — процент на менаџери кои го имаат играчот.
  \item \texttt{transfers\_in\_event} — трансфери во текот на тековниот рок.
  \item \texttt{transfers\_out\_event} — трансфери надвор во текот на тековниот рок.
  \item \texttt{status} — статус на достапност (a=достапен, d=сомнителен, i=повреден, u=неизвесен, s=суспендиран).
  \item \texttt{fixture\_difficulty\_rating} (FDR) — официјалната тежина на следниот натпревар (1--5).
\end{itemize}

\subsection{Комплетна листа на карактеристики}

Карактеристиките се организирани во шест групи:

\textbf{Тркалачки форм на играч}: \texttt{form\_last3}, \texttt{form\_last5}, \texttt{avg\_points\_per\_game\_season}, \texttt{goals\_per\_game\_season}, \texttt{assists\_per\_game\_season}, \texttt{clean\_sheet\_rate\_season}, \texttt{minutes\_reliability\_season}, \texttt{points\_per\_million}.

\textbf{Претходна лига (за дебитанти)}: \texttt{has\_prev\_league\_data}, \texttt{prev\_adjG\_per\_90}, \texttt{prev\_adjA\_per\_90}, \texttt{prev\_league\_multiplier}, \texttt{prev\_seasons\_available}, \texttt{prev\_reliability\_avg}.

\textbf{Тимски форм}: \texttt{team\_xG\_last5}, \texttt{team\_xGA\_last5}, \texttt{team\_xG\_season\_avg}, \texttt{team\_xGA\_season\_avg}, \texttt{team\_attacking\_strength}, \texttt{team\_defensive\_strength}, \texttt{team\_cs\_probability}.

\textbf{Форм на противникот}: \texttt{opp\_xG\_last5}, \texttt{opp\_xGA\_last5}, \texttt{opp\_attacking\_strength}, \texttt{opp\_defensive\_strength}, \texttt{opp\_cs\_probability}.

\textbf{Натпреварувачки карактеристики}: \texttt{current\_gw\_fdr}, \texttt{fixture\_trajectory\_score}, \texttt{home\_advantage}.

\textbf{Пазарни сигнали}: \texttt{transfers\_in}, \texttt{transfers\_out}, \texttt{selected}, \texttt{value}.

\textbf{Специфично за голман}: \texttt{saves\_per\_game\_season}, \texttt{prev\_cs\_rate}.

\textbf{Специфично за одбранбен играч}: \texttt{tackles\_won\_per\_90}, \texttt{interceptions\_per\_90}.

\section{LightGBM Модели за предикција}

\subsection{Зошто четири одделни модели}

Системот за точковање на FPL е фундаментално различен по позиција: голмани заработуваат поени за спасувања и чист лист; одбранбени играчи за чист лист и понекогаш за голови; средни играчи за голови, асистенции и чист лист; напаѓачи претежно за голови. Обединувањето на сите позиции во еден модел би резултирало во модел кој мора да научи четири сосема различни дистрибуции на излезна вредност истовремено, при значително поголема варијанса и помала предиктивна моќ. Поради тоа, за секоја позиција е тренирана посебна инстанца на LightGBM со соодветни карактеристики.

\subsection{Опис на алгоритмот}

LightGBM \cite{ke2017lightgbm} гради ансамбл на дрва на одлука на следниот начин. Почнувајќи од константна предикција, секое следно дрво $f_t$ е тренирано да ги предвиди псевдо-грешките на тековниот ансамбл:

\begin{equation}
r_i^{(t)} = -\left[\frac{\partial \mathcal{L}(y_i, \hat{y}_i)}{\partial \hat{y}_i}\right]_{\hat{y}_i = \hat{y}_i^{(t-1)}}
\end{equation}

каде $\mathcal{L}$ е функцијата за загуба (во нашиот случај, средна квадратна грешка), $y_i$ е вистинската вредност, и $\hat{y}_i^{(t-1)}$ е предикцијата по $t-1$ дрва. Новото дрво се фита на $r_i^{(t)}$ и се додава на ансамблот со стапка на учење $\eta$:

\begin{equation}
\hat{y}_i^{(t)} = \hat{y}_i^{(t-1)} + \eta \cdot f_t(x_i)
\end{equation}

Алгоритмот продолжува да додава дрва сè до претходно дефиниран број ($T$ = \texttt{n\_estimators}). Leaf-wise стратегијата за раст ги избира листовите со максимален придобивок:

\begin{equation}
\text{leaf}^* = \arg\max_{\text{leaf}} \text{Gain}(\text{leaf}) = \arg\max_{\text{leaf}} \frac{G_\ell^2}{H_\ell + \lambda} + \frac{G_r^2}{H_r + \lambda} - \frac{(G_\ell + G_r)^2}{H_\ell + H_r + \lambda}
\end{equation}

каде $G$ и $H$ се сумите на градиентите и хесианите за левиот ($\ell$) и десниот ($r$) дел, а $\lambda$ е регуларизациониот параметар.

\subsection{Walk-forward кросвалидација}

Временскиот карактер на FPL податоците налага временски-аware стратегија за валидација. Наместо стандардна случајна $k$-fold кросвалидација, применена е прошетна-напред (walk-forward) стратегија со пет набори:

\begin{table}[h]
\centering
\caption{Walk-forward кросвалидациска шема}
\label{tab:walkforward}
\begin{tabular}{clcl}
\toprule
\textbf{Набор} & \textbf{Тренирачки сет} & \textbf{Валидациски сет} & \textbf{Пондер} \\
\midrule
1 & 2019--20                    & 2020--21 & 1.0 \\
2 & 2019--20 до 2020--21        & 2021--22 & 1.5 \\
3 & 2019--20 до 2021--22        & 2022--23 & 2.0 \\
4 & 2019--20 до 2022--23        & 2023--24 & 2.5 \\
5 & 2019--20 до 2023--24        & 2024--25 & 3.0 \\
\bottomrule
\end{tabular}
\end{table}

Пондерирањето (1.0, 1.5, 2.0, 2.5, 3.0) ги дава поголемо значење на поновите сезони при пресметувањето на просечната MAE, бидејќи перформансите на поновите сезони се порелевантни за предвидувањето на идните.

\subsection{Онлајн преобучување}

По завршувањето на секоја натпреварувачка недела, вистинските поени на играчите се прикачуваат на тренирачкиот сет, и сите четири модели се преобучуваат целосно (full retrain) од нула. Ова е свесна инженерска одлука наспроти инкременталното учење: целосното преобучување обезбедува конзистентни резултати, избегнува акумулирање на грешки кои можат да се јават кај инкременталните ажурирања, и е остварливо во рамките на расположливото пресметковно време бидејќи датасетот е релативно мал.

\section{ILP Оптимизатор}

\subsection{Формулација на проблемот}

Проблемот на избор на тим е формулиран како ILP. Нека $x_i \in \{0, 1\}$ означува дали играчот $i$ е избран во составот, $c_i \in \{0, 1\}$ дали е капитен, и $v_i \in \{0, 1\}$ дали е вицекапитен.

\textbf{Функција на цел} (максимизирање):
\begin{equation}
\max \sum_{i} \hat{p}_i \cdot x_i + \hat{p}_{c^*} \cdot (1 + \text{cap\_mult})
\end{equation}

каде $\hat{p}_i$ е скорегираната предикција за играчот $i$, а $\text{cap\_mult}$ е 1.0 за нормален капитен или 2.0 за Triple Captain.

\textbf{Ограничувања}:
\begin{align}
\sum_i x_i &= 15 \quad \text{(вкупен состав)} \\
\sum_{i \in \text{GK}} x_i &= 2 \quad \text{(голмани)} \\
\sum_{i \in \text{DEF}} x_i &= 5 \quad \text{(одбранбени)} \\
\sum_{i \in \text{MID}} x_i &= 5 \quad \text{(средни)} \\
\sum_{i \in \text{FWD}} x_i &= 3 \quad \text{(напаѓачи)} \\
\sum_i \text{price}_i \cdot x_i &\leq 100.0 \quad \text{(буџет во милиони \pounds)} \\
\sum_{i \in \text{club}_j} x_i &\leq 3 \quad \forall j \in \text{clubs} \quad \text{(клуб лимит)} \\
\sum_i c_i &= 1 \quad \text{(точно еден капитен)} \\
\sum_i v_i &= 1 \quad \text{(точно еден вицекапитен)} \\
c_i + v_i &\leq x_i \quad \forall i \quad \text{(само изабрани играчи)}.
\end{align}

\subsection{Управување со чипови}

Системот автоматски активира специјални FPL чипови врз основа на прагови:

\begin{itemize}
  \item \textbf{Triple Captain (TC)}: Активира ако формот на најдобриот кандидат за капитен надминува \texttt{TC\_THRESH = 6.17} и \texttt{TC\_FORM\_MIN = 6.0}. Само два TC чипа можат да се активираат (вториот не порано од GW20).
  \item \textbf{Bench Boost (BB)}: Активира ако просечната предикција на бенчот надминува \texttt{BB\_THRESH = 9.0} (не порано од GW8).
  \item \textbf{Wildcard (WC)}: Активира ако повеќе од \texttt{WC\_THRESH = 5} играчи во тимот паднат под позицискиот просек.
  \item \textbf{Free Hit (FH)}: Активира на двојни натпреварувачки недели (DGW) за максимизирање.
\end{itemize}

\subsection{Стратегија на трансфери}

Системот банкира слободни трансфери (максимум 5, според правилата на FPL за 2025--26). Доколку ILP реши дека е потребно повеќе трансфери отколку банкираните слободни, системот бројот на дополнителни трансфери го одбива од резултатот по 4 поени за секој.

\section{Intel Pipeline — Разузнавачка мрежа во реално време}

\subsection{Intel 01: FPL Live Snapshot}

Модулот \texttt{intel\_01} ги превзема тековните статуси на играчите од FPL API. Секој играч добива бодување врз основа на официјалниот статус: достапен (100\%), сомнителен (70\%), повреден или суспендиран (0\%). Овие вредности служат како почетна точка за интегрираниот скор на достапност.

\subsection{Intel 02: Скрапирање на прес-конференции}

Модулот \texttt{intel\_02} ги скрапира FPL-ориентираните написи за прес-конференции од Fantasy Football Scout (FFS). За секој клуб наведен во насловот на написот, текстот се анализира за клучни зборови поврзани со повреди (\textit{doubt}, \textit{injury}, \textit{out}, \textit{fitness}, \textit{concern}) и достапност (\textit{available}, \textit{fit}, \textit{fine}, \textit{ready}). Резултатот е оценка на прес-конференцијата по играч (0--100).

Позната ограниченост: клубови кои не се спомнати во насловот на написот (на пример, Newcastle United во некои недели) не се анализираат, па нивните повредени играчи може да останат неоткриени.

\subsection{Intel 03: Скор на достапност}

Модулот \texttt{intel\_03} ги комбинира оценките од Intel 01 и Intel 02 во интегриран скор на достапност за секој играч:

\begin{equation}
\text{availability\_pct} = 0.65 \cdot \text{press\_score} + 0.35 \cdot \text{fpl\_score} + 5 \cdot \mathbb{1}[\text{двата извора се согласуваат}]
\end{equation}

Тежината 65\% за прес-конференциите рефлектира дека тренерите честопати дават поконкретни информации на прес-конференциите отколку FPL API, кој само ги рефлектира официјалните статуси.

\subsection{Intel 04: Скор на ризик за ротација}

Модулот \texttt{intel\_04} го квантифицира ризикот дека играчот нема да биде стартер за следната натпреварувачка недела. Скорот се пресметува врз основа на пет сигнали:

\begin{enumerate}
  \item \textbf{Стапка на старт} (\texttt{start\_rate}): Процент на натпревари во кои играчот бил во почетен состав.
  \item \textbf{Волатилност на минути} (\texttt{minutes\_volatility}): Стандардна девијација на одиграните минути.
  \item \textbf{Стапка на бенч} (\texttt{bench\_rate}): Процент на натпревари кои играчот ги почнал на клупа.
  \item \textbf{Последен тренд} (\texttt{recent\_trend}): Дали минутите на играчот се намалуваат во последните три натпревари.
  \item \textbf{Клучни зборови за ротација на прес} (\texttt{press\_keywords}): Споменување на зборови поврзани со ротација на прес-конференции.
\end{enumerate}

Финалниот скор е нормализирана тежинска сума на овие пет сигнали во опсегот [0, 100].

\subsection{Intel 05: LLM препораки}

Модулот \texttt{intel\_05} го повикува Gemini 2.5 Flash API со структуриран промпт кој вклучува тековниот состав, форм-табела, распоред на натпревари и intel сигналите. Моделот враќа структурирани препораки: капитенски избор, диференцирани играчи, цели за трансфер и ризично упатување. Овие препораки се нарациски по природа и служат за корисничкиот интерфејс, а не директно за оптимизаторот.

\subsection{Intel 06: Примена на казни}

Модулот \texttt{intel\_06} ги применува казните врз предикциите пред оптимизацијата. Формулата е:

\begin{equation}
\text{avail\_mult} = \frac{\text{availability\_pct}}{100}
\end{equation}

\begin{equation}
\text{rot\_mult} = \begin{cases}
0.40 & \text{ако } \text{rotation\_risk} \geq 80 \\
0.60 & \text{ако } \text{rotation\_risk} \geq 60 \\
0.80 & \text{ако } \text{rotation\_risk} \geq 40 \\
1.00 & \text{инаку}
\end{cases}
\end{equation}

\begin{equation}
\hat{p}_{\text{adjusted}} = \hat{p}_{\text{original}} \times \text{avail\_mult} \times \text{rot\_mult}
\end{equation}

Скоригираните предикции $\hat{p}_{\text{adjusted}}$ се проследуваат до ILP оптимизаторот наместо оригиналните предикции.

\subsection{Intel 07: Bенч Lookahead}

Модулот \texttt{intel\_07} имплементира lookahead анализа за натпреварувачките недели пред очекуваниот Bench Boost. Ако системот предвидува дека ќе биде активиран Bench Boost за N недели, играчите на бенчот добиваат бонус во нивните предикции, поттикнувајќи ILP да избере посилен бенч.

\section{Оптимизација на хиперпараметрите}

\subsection{Joint Random Search (250 обиди)}

Во првата фаза на оптимизација, беше спроведено joint случајно пребарување над целиот простор на конфигурации: ML хиперпараметри (тип на модел, длабочина на дрво, стапка на учење, итн.) заедно со симулаторските параметри (TC\_THRESH, BB\_THRESH, FDR\_MULT, итн.). Вкупно беа извршени 250 обиди, при секој обид целосна симулација на GW1--28 (вкупно ~83\,000 ILP решавачки повици).

Просторот за пребарување вклучуваше: тип на модел (XGBoost или LightGBM), n\_estimators (100--500), max\_depth (2--7), learning\_rate (0.01--0.2), subsample (0.6--1.0), colsample\_bytree (0.4--1.0), и 12 симулаторски параметри.

Клучниот наод е дека \textbf{LightGBM доминира}: 14 од 15-те најдобри резултати се постигнати со LightGBM (наспроти само еден со XGBoost). Најдобриот LightGBM обид беше Trial 220 со 1760 поени.

\subsection{Optuna Bayesian оптимизација (100 обиди)}

Во втората фаза, беше употребена Bayesian оптимизација преку Optuna \cite{akiba2019optuna} со TPE семплер. За топол старт, Trial 1 беше засеан со параметрите на Trial 220 од случајното пребарување. Ова обезбеди дека Bayesian пребарувањето почнува од добро позната регија на хиперпараметарскиот простор и ги истражува соседните конфигурации.

По 100 обиди, Trial 429 постигна \textbf{1799 поени} — подобрување за 39 поени над Trial 220. Финалните хиперпараметри на Trial 429:

\begin{table}[h]
\centering
\caption{Финални хиперпараметри на системот (Optuna Trial 429)}
\label{tab:hyperparams}
\begin{tabular}{llc}
\toprule
\textbf{Параметар} & \textbf{Опис} & \textbf{Вредност} \\
\midrule
\texttt{MODEL\_TYPE}        & Алгоритам               & lgbm   \\
\texttt{n\_estimators}      & Број на дрва            & 300    \\
\texttt{max\_depth}         & Максимална длабочина    & 3      \\
\texttt{learning\_rate}     & Стапка на учење         & 0.03   \\
\texttt{subsample}          & Подвзорокување на редови & 0.9   \\
\texttt{colsample\_bytree}  & Подвзорокување на колони & 0.6   \\
\texttt{num\_leaves}        & Максимум листови        & 31     \\
\texttt{min\_child\_samples}& Минимум примери по лист & 30     \\
\midrule
\texttt{FDR\_MULT}          & FDR-пондер (MID/FWD)    & 0.028  \\
\texttt{FDR\_MULT\_DEF}     & FDR-пондер (GK/DEF)     & 0.084  \\
\texttt{OWN\_BOOST\_GW1}    & Бонус за сопственост GW1 & 0.213 \\
\texttt{TC\_THRESH}         & Праг за Triple Captain   & 6.17  \\
\texttt{BB\_THRESH}         & Праг за Bench Boost      & 9.0   \\
\texttt{DGW\_PRED\_MULT}    & Множител за DGW         & 2.0    \\
\texttt{BENCH\_BONUS\_NORMAL} & Бонус за бенч       & 2.71   \\
\texttt{CAP\_FORM\_GATE}    & Минимален форм за капитен & 6.57 \\
\texttt{WC\_THRESH}         & Праг за Wildcard         & 5     \\
\bottomrule
\end{tabular}
\end{table}

\subsection{Споредба на методите на пребарување}

Bayesian оптимизацијата (Optuna TPE) постигна подобрување за 39 поени со само 100 обиди, наспроти 250 обиди кај случајното пребарување. Ова потврдува дека TPE ефикасно ги насочува пребарувањата кон регии со висока веројатност за добри резултати, особено при топол старт.

\section{Stage 9: LLM Нарациски слој}

После симулацијата, Stage 9 го повикува Claude API (модел \texttt{claude-sonnet-4-20250514}) за секоја натпреварувачка недела со цел генерирање на нарациско објаснување за секоја одлука. Промптот вклучува:

\begin{itemize}
  \item Избраниот тим и промените наспроти претходниот тим.
  \item Предикциите и формот на секој избран играч.
  \item Соодветниот натпревар и FDR.
  \item Активираните чипови и причините за нив.
\end{itemize}

Излезот е JSON со поле \texttt{narrative} за секое натпреварувачко коло, складиран во \texttt{models/stage9\_explanations.json}. Овој слој служи за интерпретабилност — корисникот може да разбере зошто системот донел одредена одлука — но не влијае директно врз оптимизациониот процес.

\begin{lstlisting}[language=Python, caption=Пример повик на Claude API во Stage 9]
import anthropic

client = anthropic.Anthropic(api_key=api_key)
message = client.messages.create(
    model="claude-sonnet-4-20250514",
    max_tokens=1200,
    temperature=0,
    messages=[{"role": "user", "content": prompt}]
)
narrative = message.content[0].text
\end{lstlisting}
